%==============================================================================
% MATH 5080
% Class:        9
% Date:         Sept 22nd
%==============================================================================

\newpage

\section{Limits and Continuity of Functions}

\subsection{Limits of Functions}

\begin{definition}[Limits of functions]
    Let $f: D \to \mathbb{R}, D \subset \mathbb{R}$,  and $c$
    be an accumulation point of $D$. We say $L \in \mathbb{R}$
    is the limit of $f$ at $c$ if $\forall \ep > 0, \exists \delta > 0$
    s.t. $|f(x) - L | < \ep$ whenever $0 < |x - c| < \delta$ and 
    $x \in D$,
    denoted 
    \[
        \lim_{x \to c} f(x) = L
    \]
    
\end{definition}

\begin{remark}
    \hfill
    \begin{enumerate}
        \item We define limits of functions only on accumulation points
        in the domain. 
        \item Note that $0 < |x - c|$, i.e. we care about $x \in N^*(c, \delta)$.
        We do not need $f(c)$ to be defined. 
        When it is defined, $\lim_{x \to c}f(x)$ has nothing to do with
        $f(c)$.
        \item The limit of a function at $c$ is unique.
    \end{enumerate}
\end{remark}

\begin{example}
    Show $\lim_{x \to 2}x^2 = 4$.
    \remind{Add in later}.
\end{example}


\begin{theorem}[Equivalence between limit of function and limit of a sequence]
    \hfill

    Let $f: D \to \mathbb{R}, c \in D^{\prime}$, TFAE
    \begin{enumerate}
        \item $\lim_{x \to c}f(x)= L$
        \item For every $ \{ a_n \} \subset D$ s.t. $a_n \to c$ but 
        $a_n \neq c, \forall n$, we have 
    \[
        \lim_{n \to \infty} f(a_n) = L
    \]
    \end{enumerate}
\end{theorem}

\begin{remark}
    \hfill
    \begin{itemize}
        \item 
    This result allows us to work with limits of sequences whenever
    we want to talk about limits of continuous functions
        \item We can make the above theorem weaker by only requiring $a_n \neq c$
        $\forall n > N$ for some $N$.
    \end{itemize} 
\end{remark}

\begin{example}
    $\lim_{x \to \infty} \sin \left( \frac{1}{x} \right) $ does not 
    have a limit.
\end{example}

\begin{proof}
    To show (1) $\implies$ (2): 
    \medskip
    Assume $\lim_{x \to c}f(x) = L$, 
    Fix $\{ a_n \} \subset D$ s.t. $a_n \to c, a_n \neq c, \forall n \in \mathbb{N}$
    WTS $\lim_{n \to \infty}f(a_n) = L$. 
    \medskip

    By definition, $\lim_{x \to c}f(x) = L \implies \exists \delta$ s.t. 
    $0 < |x-c| < \delta \implies |f(x) - L | < \ep$.
    Since $a_n \to c \iff \lim_{n \to \infty} a_n = c$ and $a_n \neq c, \forall n$,
    $\exists N > 0$ s.t. $\forall n > N$
    \begin{align*}
        0 < |a_n - c | < \delta \implies |f(a_n) - L | < \delta
    \end{align*}
    \medskip

    To show (2) $\implies$ (1): Assume $f(x)$ does not converge to $L$
    as $x \to c$, 
    $\exists \ep > 0$ s.t. $\forall \delta > 0$, 
    $\exists x \in N^*(c, \delta) \cap D$ s.t. $|f(x) - L | \geq \ep$.
    \begin{itemize}
        \item Let $\delta_1 = 1, a_1 \in N^*(c, \delta_1) \cap D$ s.t. $|f(a_1) - L | \geq \ep$
        \item Let $\delta_2 = \frac{1}{2}, a_2 \in N^*(c, \delta_2) \cap D$ s.t. $|f(a_2) - L | \geq \ep$
        \item $\vdots$
        \item $\delta_n = \frac{1}{n}, a_n \in N^*(c, \delta_n) \cap D$ s.t. $|f(a_n) - L | \geq \ep$
    \end{itemize} 
    Since $a_n$ satisfies $a_n \to c, a_n \neq c$, by assumption, 
    $\lim_{n \to \infty} f(a_n) = L$. This contradicts $|f(a_n) - L | \geq \ep$.
\end{proof}

\begin{corollary}
    Let $f: D \to \mathbb{R}, c \in D^{\prime}$, TFAE
    \begin{enumerate}
        \item $f$ doesn't have a limit at $c$
        \item There exists a sequence $\{ a_n \}  \subset D$
        that converges to $c$ with $a_n \neq c, \forall n$ s.t. 
        $\{ f(a_n) \} $ diverges.
    \end{enumerate} 
\end{corollary}

\begin{remark}
This is very useful for proving that a limit DNE.
\end{remark}

\begin{proposition}[Algebraic operations of limits of functions]
    Let $f: D \to \mathbb{R}, c \in D^{\prime}$. 
    Suppose $\lim_{x \to c}f(x) = L$,
    $\lim_{x \to c}g(x) = M$. 
    Then 
    \begin{enumerate}
        \item $\lim_{x \to c}(f+ g) (x) = L + M$.
        \item $\lim_{x \to c}(f\cdot g) (x) = LM$.
        \item $\lim_{x \to c}kf(x) = kL, \forall k \in \mathbb{R}$.
        \item If $g(x) \neq 0, \forall x \in D$, and $M \neq 0$, then 
        \[
            \lim_{x \to c} \frac{f}{g}(x) = \frac{L}{M}
        \]
    \end{enumerate} 
\end{proposition}

\begin{definition}[One sided limits]
    \hfill
    \begin{enumerate}
        \item We say that $\lim_{x \to c^-}f(x) = L$ if 
        $\forall \ep > 0$, $\exists \delta > 0$ s.t. 
        $c - \delta < x < c \implies |f(x) - L | < \ep$ 
        \item We say that $\lim_{x \to c^+}f(x) = L$ if 
        $\forall \ep > 0$, $\exists \delta > 0$ s.t. 
        $c < x < c  + \delta\implies |f(x) - L | < \ep$ 
    \end{enumerate} 
\end{definition}

\begin{proposition}[
        Equivalence of one-sided and two-sided limits
    ]
    \begin{align*}
        \lim_{x \to c}f(x) = L
        \iff
        \lim_{x \to c^-}f(x) =
        \lim_{x \to c^+}f(x) =
        L
    \end{align*}
\end{proposition}

\subsection{Continuous functions}


\begin{definition}[Continuity]
    Let $f: D \to \mathbb{R}$
    \begin{enumerate}
        \item (\textbf{Continuity at a point})
    Let $c \in D$,  
    we say $f$  is \textbf{continuous at $c$} if 
    $\forall \ep > 0, \exists \delta > 0$ s.t. 
    $|f(x) - f(c)| < \ep$ whenever 
    $|x - c| < \delta$.
        \item (\textbf{Continutiy over a set}): Let $S \subset D$
    if $\forall c \in S$, $f$ is continuous at $c$, then we say 
    that $f$ is \textbf{contiuous on $S$}.
        \item (\textbf{Continuous function}): if $f$ is continuous
    on $D$, then we say $f$ is a \textbf{continuous function}.
    \end{enumerate} 
\end{definition}

\begin{remark}
    \hfill
    \begin{enumerate}
        \item Compared to the definition of 
        the limit of a function at a point,
        we now require $c \in D$ instead of $c \in D^{\prime}$.
        \item We now allow $x = c$.
        \item We cannot immediately apply the previous definition 
        of limit of a function at point to define continuity
        because 
        we don't know that $c$ is an accumulation point.
        \item Depending on whether $c$ is an accumulation point
        or an isolated point, we have different conditions
        that are required for the continuity of $f$ at $c$. 
        Refer to the following two propositions.
        % \item If $c$ is an isolated pint of the domain, then any function is 
        % continuous at $c$.
        % \item If $c$ is an accumulation point, $f$ continuous at $c$ $\iff$
        % $\lim_{x \to c}f(x) = f(c)$.
    \end{enumerate} 
\end{remark}

\begin{proposition}[Any function is continuous at isolated point]
    If $c$ is an isolated point of $D$, then 
    $f$ is always continuous at $c$.
\end{proposition}

\begin{proof}
    $c$ isolated $\implies \exists \delta_0$ s.t.
    $N^*(c, \delta_0) \cap D = \emptyset$. 
    Fix any $\ep > 0$, pick $\delta = \delta_0$. 
    Then, 
    \begin{align*}
        x \in N(c, \delta) \implies x = c \implies
        |f(x) - f(c) | = 0 < \ep
    \end{align*}
\end{proof}


\begin{proposition}
    If $c$ is an accumulation point of $D$, then 
    $f$ is continuous at $c$ if and only if $\lim_{x \to c}f(x)=$
    $f(c)$.
\end{proposition}

\begin{proof}
    By definition, \remind{exercise}.
\end{proof}




\begin{proposition}[Equivalent characterization of continuity]
    $f$ is continuous at $c$ if any only if for any neighborhood 
    of $f(c)$, say $V = N(f(c), \ep), \forall \ep$, 
    there exists a neighborhood $U = N(c, \delta)$ such that 
    $f(U \cap D) = V$.
\end{proposition}

\begin{theorem}
    Let $f: D \to \mathbb{R}, c \in D$ , TFAE
    \begin{enumerate}
        \item $f$ is continuous at $c$
        \item If $\{ a_n \} \subset D $ s.t. $a_n \to c$, then 
        \[
            \lim_{n \to \infty} f(a_n) = f(c)
        \]
    \end{enumerate} 
\end{theorem}

\begin{proof}
    \remind{revisit this proof}
    \medskip

    $c$ is either an isolated point or an accumulation point. 
    \begin{enumerate}
        \item Case 1: $c$ isolated.
        We want to show that (1) and (2) are always true. 
        \begin{itemize}
            \item We already know that (1) is always true, 
            any $f$ is continuous at $c$ if $c$ isolated.
            \item WTS (2). Fix any $\{ a_n \}$ converging to $c$. 
            Since $c$ isolated, $\exists \ep > 0$ s.t. 
            $N^*(c, \ep) \cap D = \emptyset$. 
            Since $a_n \to c, \exists N$ s.t. $\forall n > N$,
            $|a_n - c | < \ep \implies a_n = c \implies f(a_n) = f(c)$.
        \end{itemize} 
        \item Case 2: $c$ accumulation point
            \begin{itemize}
                \item (1) $\implies$ (2): 
                \begin{align*}
                    f \text{continuous at } c 
                    & \implies \lim_{x \to c}f(x) = f(c) \text{ by definition} \\
                    & \implies \{ a_n \} \subset D, a_n \neq c, a_n \to c \implies 
                    \lim_{n \to \infty} f(a_n) = f(c)
                \end{align*}
                ATC otherwise, i.e. $\{ a_n \} \subset D, a_n \to c, \lim_{n \to \infty}f(a_n) \neq f(c)$
                then $\exists \ep > 0$ and some subsequence $\{ a_{n_k} \} $
                such that $a_{n_k} \to c$ but $ \left| f(a_{n_k}) - f(c)\right|  > \ep$.
                Hence $a_{n_k} \neq c, \forall k$. However, this contradicts with
                the result.
                \item (2) $\implies$ (1): 
                Assume (2), then we immediately have 
                \begin{align*}
                    &\{ a_n \} \subset D, a_n \neq c, a_n \to c \implies 
                    \lim_{n \to \infty} f(a_n) = f(c) \\
                    & \implies \lim_{x \to c}f(x) = f(c) \\
                    & \implies f \text{ continuous at } c
                \end{align*}
            \end{itemize} 
    \end{enumerate} 
\end{proof}

\begin{theorem}
    $f, g: D \to \mathbb{R}, c \in D$. Suppose $f, g$ both
    continuous at $c$. Then the following functions are also
    continuous at $c$:
    \begin{enumerate}
        \item $f + g$
        \item $fg$
        \item $kf$
        \item $ \frac{f}{g}$, if $g(c) \neq 0$
    \end{enumerate} 
\end{theorem}

\begin{remark}
    The assumption that $\exists \ep > 0$
    s.t. 
    $\forall x \in N(c, \ep), g(x) \neq 0$
    is necessary and sufficient condition for $g(c) \neq 0$.
\end{remark}

\begin{proof}
    We prove using the previous theorem, and the 
    analogous result for limits.
\end{proof}


\begin{theorem}
    Let $f: D \to \mathbb{R}, g: E \to \mathbb{R}$ s.t.
    $f(D) \subset E$, 
    suppose $f$ continuous at $c \in D$, $g$ continuous at $f(c) \in E$, 
    then $g \circ f : D \to \mathbb{R}$ is continuous at $c$.
\end{theorem}

\begin{proof}
    Fix a neighborhood $V$ of $g \circ f(c)$, we seek a neighborhood $U$
    such that $g \circ f( U \cap D) \subset V$. 
    \medskip

    Since $g$ continuous at $f(c)$, $\exists $ a neighborhood $W$
    of $f(c)$ such that $g(W \cap E) \subset V$. 
    \medskip

    Since $f$ continuous at $c$, $\exists$ a neoighboord $U$ of $c$ such that 
    $f(U \cap D) \subset W$.
\end{proof}

\begin{remark}
    $f(D) \subset E \implies f(U \cap D) \subset W \cap E$
    $\implies g \circ f(U \cap D) \subset g(W \cap E) \subset V$.
    \begin{center}
    \begin{tikzpicture}[scale=0.9, every node/.style={font=\small}]
        % domain D and neighborhood U of c
        \draw[thick] plot[smooth cycle, tension=0.8]
            coordinates {(-1.5,0.4) (-1.0,0.6) (-0.7,1.2) (0.1,0.8) (0.6,1.1) (1.2,0.5) (0.9,-0.2) (1.3,-0.9) (0.2,-0.8) (-0.5,-1.3) (-1.1,-0.7) (-0.9,-0.1)};
        \draw[thick, blue] (0.3,0.3) circle (0.8);
        \fill (0.3,0.3) circle (1.5pt) node[below left] {$c$};
        \node at (-1.7,0.9) {$D$};
        \node[blue] at (0.3,1.35) {$U$};

        % E containing f(D), and neighborhood W of f(c)
        \begin{scope}[xshift=5.5cm]
            \draw[thick] plot[smooth cycle, tension=0.8]
                coordinates {(-1.7,0.3) (-1.3,1.1) (-0.6,1.6) (0.2,1.5) (0.7,1.8) (1.5,1.2) (1.7,0.4) (2.0,-0.4) (1.3,-1.1) (0.4,-1.0) (-0.3,-1.6) (-1.2,-1.3) (-1.4,-0.5)};
            \draw[thick, dashed, red!70!black] plot[smooth cycle, tension=0.8]
                coordinates {(-0.6,0.5) (-0.2,0.9) (0.4,0.7) (0.9,0.3) (1.5,0.0) (1.6,-0.5) (1.1,-0.8) (0.4,-0.6) (-0.3,-0.9) (-0.8,-0.4)};
            \draw[thick, blue] (0,0.3) circle (1.05);
            \draw[thick, violet] (0,0.3) ellipse (0.55 and 0.42);
            \fill (0,0.3) circle (1.5pt) node[below] {$f(c)$};
            \node[violet] (fUD) at (-1.9,-1.9) {$f(U \cap D)$};
            \draw[violet, thin] (fUD.north) -- (-0.52,0.16);
            \node[red!70!black] at (1.2,-0.4) {$f(D)$};
            \node[blue] at (-1.25,0.55) {$W$};
            \node at (0.6,-1.5) {$E$};
        \end{scope}

        % V around g(f(c))
        \begin{scope}[xshift=10.5cm]
            \draw[thick] plot[smooth cycle, tension=0.8]
                coordinates {(-1.4,0.5) (-0.9,1.0) (-0.2,0.9) (0.4,1.3) (1.2,0.8) (1.0,0.1) (1.5,-0.6) (0.6,-1.1) (-0.2,-0.7) (-0.9,-1.2) (-1.5,-0.4)};
            \fill (-0.4,0) circle (1.5pt) node[right] {$g(f(c))$};
            \node at (1.6,1.0) {$V$};
        \end{scope}

        % maps
        \draw[-{Stealth}, thick] (1.1,0.8) to[bend left=25] node[above] {$f$} (4.3,0.8);
        \draw[-{Stealth}, thick] (6.6,0.8) to[bend left=25] node[above] {$g$} (9.4,0.8);
    \end{tikzpicture}
    \end{center}
\end{remark}

\begin{theorem}[Limit passes through continuous functions]
    \hfill

    Let $f: D \to \mathbb{R}, g: E \to \mathbb{R}$,
    $f(D) \subset E$, 
    suppose $\lim_{x \to c}f(x) = L \subset E$ and 
    $g$ is continuous at $L$,
    then 
    \[
        \lim_{x \to c} g(f(x)) = 
        g \left( \lim_{x \to c} f(x)\right)  = 
        g(L)
    \]
\end{theorem}

\begin{example}
    If however, we only have 
    $\lim_{x \to c} f(x) = L$,
    $\lim_{y \to L} g(y) = M$,
    then in general we \textbf{do not have}
    $\lim_{x \to c} g(f(x)) = M$.
    \medskip

    For example, take $f(x) = L$  a constant function, 
    then $\lim_{x \to c}g(f(x))= g(L)$,
    which is not necessarily equal to 
    $\lim_{y \to L}g(y)$, unless we know
    that $g$ is continous at $L$.
\end{example}

\begin{proof}
    Fix $\ep >0$. WTS $\exists \delta > 0$ s.t. 
    if 
    $0 < |x - c | < \delta $ and $x \in D$, 
    then $|g(f(x)) - g(L) | < \ep$. 
    \medskip

    By definition of continituity of $g$ at $L$,
    for $\ep$, $\exists \eta > 0$ s.t. 
    $|y - L | < \eta, y \in E \implies $
    $|g(y) - g(L) | < \ep$
    \medskip

    By definition of limit of $f$ at $c$,
    since $\lim_{x \to c}f(x) = L$,
    $\exists \delta > 0$ s.t. 
    $0 < | x - c | < \delta, x \in D$, 
    $\implies |f(x) - L | < \eta$.
    \medskip

    Since $f(D) \subset E$, 
    $0 < | x - c | < \delta, x \in D$
    $\implies |f(x) - L | < \eta, f(x) \in E$
    $\implies |g(f(x)) - g(L) | < \ep$.
\end{proof}

\begin{example}
    $f(x) = x^2, D = (-1, 1), f(D) = [0, 1)$
\end{example}

\begin{remark}
    Continuous functions do not necessarily preserve
    openess.
\end{remark}

\begin{theorem}[
    Preimage of open set in $\mathbb{R}$
    is open
]
Let $f$ be a function on $\mathbb{R}$, TFAE
\begin{itemize}
    \item $f$ is continuous on $\mathbb{R}$.
    \item For any open $V \subset \mathbb{R}, f^{-1}(V)$ is open.
\end{itemize} 
\end{theorem}

\begin{remark}
    This can be made stronger into the following theorem.
\end{remark}

\begin{theorem}[Preimage of open is open]
    Let $f: D \to \mathbb{R}$, TFAE 
    \begin{enumerate}
        \item $f$ is continuous on $D$
        \item Any open image has a preimage that is the 
        intersection of an open set and the domain.
        \[
        \forall V \subset \mathbb{R}, 
        V \text{ open},
        \exists U \subset \mathbb{R} 
        \text{ s.t. }
        f(U \cap D) = V, 
        \text{ i.e. }
        f^{-1} = U \cap D.
        \]
    \end{enumerate} 
\end{theorem}

\begin{remark}
    We want to say that any open image has an 
    open pre-image, but that is not always true.
    We can only say $f^{-1}(V) = U \cap D$
    for some open $U$.
\end{remark}

\begin{corollary}
    Let $f: D \to \mathbb{R}$ , TFAE
    \begin{enumerate}
        \item $f$ is continuous on $D$
        \item for any closed $V \subset \mathbb{R}$,
        $\exists U \subset \mathbb{R}$, $U$ closed, 
        s.t. 
        $f^{-1}(V) = D \cap U$.
    \end{enumerate} 
\end{corollary}

\begin{proof}
    We prove a special case where $D = \mathbb{R}$.
    \medskip

    (1) $\implies$ (2): Assume $f$ continous on $D$, let $V$ be open. WTS $f^{-1}(V)$
    open. 
    \medskip

    Fix any $c \in f^{-1}(V)$. WTS $\exists \ep > 0$  s.t $N(c, \ep) \subset f^{-1}(V)$.
    \medskip

    Since $c \in f^{-1}(V), f(c) \in V$. Since $V$ open, $\exists$ a neighborhood
    of $f(c)$ say $U$ s.t. $U \subset V$.
    \medskip
    
    Since $f$ continuous at $c$, by definition, $\exists \ep > 0$  s.t. 
    $f(N(c, \ep)) \subset U \subset V$. Hence $N(c, \ep) \subset f^{-1}(V)$.
    \medskip

    (2) $\implies$ (1): WTS $f$ continuous at any $c \in \mathbb{R}$. Let $V$ 
    be any neighborhood of $f(c)$. WTS $\exists$ a neighborhood of $c$, say $U$ s.t.
    $f(U) \subset V$. 
    By assumption, we know $f^{-1}(V)$ is open. Since $c \in f^{-1}(V)$, 
    $c$ must be an interior point of $f^{-1}(V)$. So there exists a neighborhood 
    of $c$ say $U$ s.t. $U\subset f^{-1}(V)$, i.e. $f(U) \subset V$.
\end{proof}













